%%-----------Chapters start-------------------------------------
%%-----------Chapter 1------------------------------------------
\chapter{Introduction}\label{chap:intro}
Microbes by themselves can produce chemical molecules that have significant use as industrial chemicals or their precursors \cite{Handelsman07, Fuhrman09}. These single microbial species are limited in their metabolic capacity. However, the interactions between different microbial species in a microbial community can result in emergent metabolic behaviors that support the production of chemical species that can have multiple uses, e.g., aid the process of stabilizing the ecosystem, enrich agriculture \cite{Kaye05} and maintain stable human health \cite{Ley10, Greenblum12, Qin12}. 
% Hence, it is of utmost importance to evaluate the metabolic engineering potential of metabolic communities found in nature. 
% Metabolic engineering techniques have successfully been applied to design and optimize the microbial production of antibiotics, alcohols and amino acids by single strains \cite{Garcia16}. 
Presently, research efforts are beginning to focus on analysis and control of the metabolism of complex microbial communities \cite{Chiu14, Garcia16}.  %% this is important to our work, keep statement in   

 In the field of bioenergy, efficient consumption and conversion of  both hexose and pentose sugar monomers in biomass hydrolysates is an important requirement for the commercialization of second-generation biofuels \cite{Hanly14}. Biomass derived from plant matter is a heterogeneous material that yields a variety of pentose and hexose sugars when subjected to enzymatic hydrolysis \cite{Hanly11, Zaldivar01}. Microbial co-cultures/consortia have been shown to be more effective for fermentation of sugar mixtures to ethanol, compared to individual strains \cite{Lynd02, Hanly11}. Instead of attempting to engineer cells to deviate from their natural metabolic behavior, these consortia take advantage of the native capabilities of different species to metabolize different sugars and direct available carbon to targeted biofuels or biofuel precursors. This strategy allows metabolic engineering efforts to focus on leveraging natural microbe behavior instead of introducing foreign genes whose effects are often counteracted upon \cite{Hanly11, Kuyper05}. In other words, microbial mixed cultures allow the combination of the inherent conversion strengths of two or more microbes \cite{Hanly13, Brenner08}. 

%A variety of microbial co-cultures have been proposed for fermentation of lignocellulosic hydrolysates \cite{Hickert13, Yadavt13}


%Another outstanding problem in the conversion of lignocellulosic materials to commodity chemicals is the production of chemical compounds such as weak acids, furan derivatives, and phenols that can have a deleterious effect on fermentative micro-organisms \cite{Palmqvist00}. 

The development of microbial consortia for efficient biofuel production is a very challenging problem. Experimental data on the behavior of microbial consortia is limited, especially for species combinations that do not appear naturally \cite{Fazzinit10, Ramadas09, VerBerkmoes09}. Moreover, purely experimental approaches are not well suited to handle the wide variability of available biomass feedstocks and the immense diversity of microbial species. Metabolic modeling offers the potential for the identification of synergistic combinations of biofuel-producing microbes and optimization of fermentation conditions for efficient substrate utilization and high product yields, thereby informing and guiding experimental efforts. Though there have been several approaches to create metabolic models for consortia \cite{Biggs15}, none of these approaches involve a well-defined formal method based approach. In this work, a hybrid systems approach is formulated to model metabolism in microbial co-cultures. This allows the introduction of formal methods from computer science/systems engineering to the realm of consortia bio-engineering.

%The existing systems biology literature provides approaches to simulate microbial metabolism directly from the stoichiometry of the organism \cite{Chiu14}.

For a metabolic network whose stoichiometry can be defined in sufficient detail, a number of analyses and simulation methods, termed \textit{constrained-based metabolic modeling} can be brought to bear as reviewed by Chang et al \cite{Chang08}. The stoichiometry of the metabolic network encodes a great deal of information about the associated metabolism, including the network topology and allows only certain combinations of fluxes in the steady state. In the constrained-based metabolic modeling approach, a phenotypically feasible solution space is defined by specifying certain assumptions, e.g., a steady-state assumption \cite{Heirendt19}, and the solution space can be further reduced by the imposition of constraints related to physico-chemical laws, genetics, and the environment \cite{Heirendt19, O'Brien15}.  Thus, constraint-based techniques often focus on computing the possible steady states of a biochemical reaction network, given the stoichiometry of the network \cite{Wang17}. 
%Flux balance analysis (FBA) is a constraint-based method that maximizes a linear objective function formed by multiplying each reaction flux by a predetermined coefficient, subject to a steady-state assumption, as well as lower and upper bounds on each reaction flux\cite{Orth10, Heirendt19}.

%% this might be useful to add to FBA section

% If there is a nonredundant list of $m$ compounds participating in a non-redundant list of $n$ chemical reactions that make up a metabolic network, we may define an $m \cdot n$ matrix $S$, called the stoichiometric matrix, whose elements $S_{ij}$ have magnitudes equal to the stoichiometry of compound $i$ in reaction $j$. Compounds consumed in a reaction (reactants) are assigned negative stoichiometry values, those produced (products) positive stoichiometry values. Even in the absence of any information about the reaction rate laws, S encodes a great deal of information about the associated
% metabolism, including the network topology and permissible combinations of fluxes in the steady state, when compound concentrations are constant with time despite fluxes producing and consuming the compound. With sufficient physical or empirical constraints on the permissible fluxes, flux balance analysis and metabolic control analysis can elucidate much about the steady-state behaviors of metabolism. Availability of more detailed thermodynamic constraints with respect to individual reactions leads to energy balance analysis.

Most common among such analyses are methods like flux balance analysis (FBA) and its variants. Flux balance analysis (FBA) is a constraint-based method that maximizes a linear combination of reaction fluxes in the metabolic network, subject to a steady-state assumption, as well as lower and upper bounds on each reaction flux\cite{Orth10, Heirendt19}. FBA and its variants have usually been used to characterize metabolic interactions within a single microbe for applications in environmental and industrial biotechnology \cite{Garcia16}. However, this approach can be extended to co-cultures of two or more microbes as described in Salimi et al. \cite{Salimi10} and Perez-Garcia et al. \cite{Garcia16}.
%Some other constraint-based methods include metabolic control analysis, elemental mode analysis and extreme pathway analysis. 


% A co-culture metabolic model of {\it{Escherichia coli}}  and {\it{Saccharomyces cerevisiae}} strains engineered to selectively utilize only glucose and  only xylose respectively has been developed by Hanly et al. \cite{Hanly11}. Another co-culture metabolic model has been developed by the same group of researchers \cite{Hanly14} for a respiratory deficient mutant of the ethanologic yeast {\it{Saccharomyces cerevisiae}} and the xylose-fermenting wild-type yeast {\it{Scheffersomyces stipites}} to optimize fermentation strategies for the simultaneous detoxification and fermentation of lignocellulosic hydrolysates. 

%While FBA is computationally efficient and produces accurate results, it lacks a number of properties that are present in formal method-based model checking systems. These include but are not strictly limited to: traceability, explainability, exhaustive analyzability and visualizability. %Thus, we're able to (during a simulation, or afterwards) visualize a) which states were active and when, b) which transitions fired and when, and can effectively represent the entire simulation in a graph-like structure where the simulation is a unique path through it.

By building on existing approaches in constraint-based metabolic modeling, particularly as applied to a simplistic microbial co-culture, a hybrid systems modeling framework for microbial co-cultures has been developed that preserves the phenotypically-defined structure of the solution space. In section \ref{sec:pre-existing}, the lack of hybrid systems approaches to model microbial co-cultures is detailed, followed by a description of constraint-based metabolic modeling methods such as FBA and dynamic FBA in section \ref{sec:fba}. We then elaborate our approach in section \ref{sec:model} and discuss the results in section \ref{sec:res}. 

%elementary mode analysis and extreme pathway analysis 

Next, the hybrid systems modeling approach is formulated with respect to a simplified, ``toy" system -- namely, an {\it{in silico}} co-culture of two microbes with trivial metabolic networks ($<$ 10 reactions each) that exhibit emergent metabolic behavior when placed in the same medium under specific iniial conditions. This toy system in particular is considered as it has been already used to show \textit{in silico} emergent metabolite production under the dynamic FBA approach\cite{Chiu14}. Then, we discuss the model predicted behavior of the toy microbe co-culture in section \ref{sec:bio}, followed by the modeling paradigm using hybrid systems approach discussed in section \ref{sec:hybrid}. Next, we model the microbe interactions as hybrid systems discussed in section \ref{sec:res} and how our approach differs from previous approaches. Finally, we formalize the properties to check if they are satisfied by the modeled behavior. Our initial results, we see similar performance as elaborated by FBA analysis. 


% \textbf{Comment: Why do we add descriptions of anything other than FBA, if we are not comparing our results with those methods ???}

% \textbf{Comment (JFB): FBA and dFBA were directly used to build our approach (it is a reduction from what is computed by dFBA to an automata that computes the same). Extreme pathways find incorporation in our approach on larger, genome-scale models, and justify the more basic work performed here (even the basic null space analysis gives results that are qualitatively similar to what you would expect to see from expas). Elementary modes were not incorporated but were researched in the process and may find use in later application, as they also form a convex basis for the flux cone and thus can be used interchangeably with expas or any other basis. }